\chapter{MARCO TEÓRICO}


% ------------------------------------------------------------
\section{NB-IoT: arquitectura y mecanismos de cobertura extendida}
% ------------------------------------------------------------

Para hacer frente al escenario descrito, el proyecto 3GPP introdujo en su Release 13 el
Internet de las Cosas de Banda Estrecha (NB-IoT, \textit{Narrowband Internet of Things}),
una tecnología de acceso radio diseñada para soportar comunicaciones MTC masivas en espectro
licenciado~\cite{ref3, ref7}. A diferencia de tecnologías de área amplia que operan en bandas
no licenciadas ---como LoRaWAN o Sigfox---, NB-IoT ofrece garantías de calidad de servicio,
mayor seguridad y la posibilidad de aprovechar la infraestructura celular existente, lo que
reduce los costos de despliegue de forma importante~\cite{ref3}.

En cuanto a la capa física, NB-IoT opera con un ancho de banda de 180 kHz, equivalente a un
bloque de recurso en LTE (\textit{Long Term Evolution}), y usa el acceso múltiple por división
de frecuencia ortogonal (OFDMA, \textit{Orthogonal
Frequency Division Multiple Access}) en el enlace descendente y una combinación de OFDMA y
acceso múltiple por división de frecuencia de portadora única (SC-FDMA, \textit{Single-Carrier
Frequency Division Multiple Access}) en el ascendente~\cite{ref7}. Aunque este ancho de banda puede parecer limitado, en realidad es una
decisión de diseño: permite que los dispositivos funcionen con radios de baja complejidad y
bajo consumo de energía, logrando baterías con vida útil de hasta diez años sin
reemplazo. El estándar tiene tres modos de despliegue ---en banda (\textit{in-band}), en
banda de guarda (\textit{guard-band}) y autónomo (\textit{standalone})--- que permiten a los
operadores integrar NB-IoT dentro de sus redes LTE existentes o en espectro reasignado de
tecnologías anteriores~\cite{ref3, ref18}.

Una característica importante de NB-IoT, que está directamente relacionada con el problema
que se estudia en este trabajo, es su mecanismo de mejora de cobertura (CE,
\textit{Coverage Enhancement}). Según la potencia de referencia de la señal recibida (RSRP,
\textit{Reference Signal Received Power}), la estación base (eNB, \textit{evolved Node B}) clasifica a cada dispositivo en
uno de tres niveles de cobertura: CE0, CE1 y CE2~\cite{ref7}. Esta clasificación define
cuántas repeticiones de transmisión necesita el dispositivo: los que están en condiciones
difíciles ---lejos de la celda, dentro de edificios o con obstáculos--- compensan la pérdida
de señal repitiendo la transmisión varias veces. La Tabla~\ref{tab:ce_params} muestra los
umbrales de RSRP y los principales parámetros de cada nivel.

\begin{table}[h!]
\centering
\caption{Niveles de cobertura en NB-IoT y parámetros NPRACH~\cite{ref7}}
\label{tab:ce_params}
\begin{tabular}{|l|c|c|c|}
\hline
\textbf{Parámetro} & \textbf{CE0} & \textbf{CE1} & \textbf{CE2} \\
\hline
Umbral RSRP & $>$ $-$115 dBm & $-$127 a $-$115 dBm & $\leq$ $-$127 dBm \\
Repeticiones MSG1 & 1 & 8 & 32 \\
Periodicidad NPRACH & 320 ms & 640 ms & 2560 ms \\
Latencia MSG1 & 6{,}40 ms & 51{,}20 ms & 204{,}79 ms \\
\hline
\end{tabular}
\end{table}

El canal físico de acceso aleatorio de banda estrecha (NPRACH,
\textit{Narrowband Physical Random Access Channel}) es el recurso que usan los dispositivos
para iniciar la conexión. En su configuración estándar, NB-IoT tiene solo 48 preámbulos
ortogonales por oportunidad de acceso~\cite{ref7}. Este número, que ya es bajo para escenarios
de densidad moderada, se convierte en un cuello de botella crítico cuando cientos o miles de
dispositivos intentan acceder al mismo tiempo, algo común en tráfico MTC de tipo ráfaga
provocado por eventos externos. La falta de preámbulos aumenta la probabilidad de colisión,
obliga a retransmisiones y eleva el retardo de acceso, sobre todo en los niveles CE1 y CE2,
donde cada intento fallido implica esperar hasta 640 ms o 2560 ms para la siguiente
oportunidad de acceso.

Esta limitación es lo que motiva el estudio de esquemas de acceso no ortogonal sobre NB-IoT.
Si el sistema logra distinguir y recuperar dispositivos que han colisionado en el mismo
preámbulo, la capacidad del canal de acceso aumenta sin necesidad de agregar más recursos de
frecuencia. En la siguiente sección se describe el procedimiento estándar de acceso aleatorio
en NB-IoT, que es la base sobre la cual se construyen los esquemas analizados en este
trabajo.


% ------------------------------------------------------------
\section{Procedimiento de acceso aleatorio en NB-IoT}
% ------------------------------------------------------------

El acceso aleatorio es el mecanismo que usa un equipo de usuario (UE, \textit{User Equipment})
de NB-IoT para establecer su primera comunicación con el eNB antes de que se le asignen
recursos dedicados. En su forma estándar, llamada acceso aleatorio basado en concesión
(\textit{grant-based random access}), el procedimiento consiste en cinco mensajes entre el
dispositivo y la red~\cite{ref5, ref7}. La Figura~\ref{fig:ra_msc} resume ese intercambio,
que se describe a continuación.

\begin{figure}[H]
\centering
\begin{tikzpicture}[
    font=\small,
    >={Stealth[length=2mm]},
    msg/.style={->, thick},
]
    % lifelines
    \node (ue) at (0,0) {\textbf{UE}};
    \node (enb) at (8,0) {\textbf{eNB}};
    \draw[dashed] (ue) -- (0,-7);
    \draw[dashed] (enb) -- (8,-7);
    % messages
    \draw[msg] (0,-0.8) -- node[above,font=\scriptsize]{MSG1: preámbulo (NPRACH)} (8,-0.8);
    \draw[msg] (8,-2.0) -- node[above,font=\scriptsize]{MSG2: RAR (TC-RNTI, \textit{timing advance})} (0,-2.0);
    \draw[msg] (0,-3.2) -- node[above,font=\scriptsize]{MSG3: solicitud de conexión RRC (NPUSCH)} (8,-3.2);
    \draw[msg] (8,-4.4) -- node[above,font=\scriptsize]{MSG4: resolución de contención (C-RNTI)} (0,-4.4);
    \draw[msg] (0,-5.6) -- node[above,font=\scriptsize]{MSG5: confirmación (RRC Setup Complete)} (8,-5.6);
\end{tikzpicture}
\caption{Intercambio de los cinco mensajes del procedimiento de acceso aleatorio basado en
concesión en NB-IoT. El UE inicia con el preámbulo (MSG1) y el enlace queda establecido tras
la confirmación (MSG5).}
\label{fig:ra_msc}
\end{figure}

El proceso comienza cuando el UE escoge al azar
uno de los preámbulos disponibles en el NPRACH y lo transmite en la oportunidad de acceso que
corresponde a su nivel CE (MSG1). Este preámbulo le indica al eNB que un dispositivo quiere
conectarse y además permite estimar el tiempo de llegada de la señal (ToA,
\textit{Time of Arrival}) para calcular el avance de temporización necesario para sincronizar
el enlace ascendente. Al detectar el preámbulo, el eNB responde con el
MSG2 o respuesta de acceso aleatorio (RAR, \textit{Random Access Response}), que
incluye el identificador del preámbulo reconocido (RAPID, \textit{Random Access Preamble
Identifier}), un identificador temporal de radio (TC-RNTI, \textit{Temporary Cell Radio
Network Temporary Identifier}) y el avance de temporización calculado. El UE que reconoce su
RAPID en el RAR transmite entonces el mensaje de capa 3 inicial (MSG3) sobre el
canal físico compartido del enlace ascendente de banda estrecha (NPUSCH, \textit{Narrowband
Physical Uplink Shared Channel}), usando el TC-RNTI asignado. El eNB responde con un
mensaje de resolución de contención (MSG4) que contiene la identidad permanente del
dispositivo (IMSI, \textit{International Mobile Subscriber Identity}) y el identificador de
radio definitivo (C-RNTI, \textit{Cell Radio Network Temporary Identifier}), con lo cual queda
establecido el enlace de control de recursos de radio (RRC, \textit{Radio Resource Control}).
Por último, el UE confirma con un
acuse de recibo (MSG5), y a partir de ese momento la conexión queda
establecida~\cite{ref7}.

Ahora bien, el punto crítico de este flujo está en el MSG1. Como los dispositivos eligen el
preámbulo de forma independiente y al azar, dos o más UEs pueden escoger el mismo preámbulo
en la misma oportunidad de acceso, lo que produce una colisión. En el acceso aleatorio
ortogonal convencional de NB-IoT (NB-ORA, \textit{NB-IoT Orthogonal Random Access}), el
eNB no puede distinguir entre los dispositivos que colisionaron y genera un solo
RAR para ese preámbulo. Todos los UEs involucrados reciben el mismo TC-RNTI y transmiten su
MSG3 sobre los mismos recursos, lo cual genera otra colisión en el canal ascendente. Solo el
mensaje MSG4, mediante la coincidencia del IMSI, permite que cada dispositivo sepa si la
resolución de contención le corresponde; los demás UEs, al no reconocer su identidad en MSG4,
deben esperar a que expire el temporizador T300 y empezar de nuevo desde MSG1~\cite{ref4,
ref7}. En los niveles CE1 y CE2, donde las oportunidades de acceso pueden estar separadas por
640 ms o 2560 ms, respectivamente, cada reintento fallido agrega retardos que pueden pasar de
varios segundos por dispositivo.

La probabilidad de colisión crece rápidamente con la densidad de dispositivos. Si en una
oportunidad de acceso hay $M$ UEs activos y $N$ preámbulos disponibles, la probabilidad de que
al menos dos dispositivos seleccionen el mismo preámbulo se puede aproximar mediante la
Ecuación~\ref{eq:pcol}~\cite{ref4}.

\begin{equation}
P_{\text{col}} = 1 - \frac{N!}{N^{M}(N-M)!}
\label{eq:pcol}
\end{equation}

Para NB-IoT, con $N = 48$ preámbulos, basta con que $M$ supere unas pocas decenas de
dispositivos para que la probabilidad de colisión se vuelva dominante, haciendo que la mayor
parte de los intentos de acceso fallen en el primer mensaje~\cite{ref7}.

Ante esta limitación, en la literatura se han propuesto dos tipos de solución. La primera usa
técnicas de control de sobrecarga, como el bloqueo de clase de acceso (ACB,
\textit{Access Class Barring}), que limitan cuántos dispositivos pueden intentar acceder en
cada intervalo, reduciendo la congestión pero aumentando el retardo~\cite{ref5}. La segunda,
y la más importante para este trabajo, propone modificar el procedimiento de acceso para
recuperar los dispositivos que colisionaron en vez de descartarlos, aprovechando las
diferencias entre sus señales para separarlos y permitir que todos avancen en el proceso de
conexión. Este enfoque se conoce como acceso aleatorio no ortogonal, y se presenta en la
siguiente sección.


% ------------------------------------------------------------
\section{Acceso Aleatorio No Ortogonal (NORA)}
% ------------------------------------------------------------

El acceso aleatorio no ortogonal surge como respuesta directa al problema de colisiones
descrito en la sección anterior. Mientras que el esquema ortogonal convencional descarta los
dispositivos que colisionaron y los obliga a reintentar, NORA propone recuperarlos: si las
señales que se superponen en el receptor tienen diferencias suficientes ---en potencia
recibida, en tiempo de llegada o en código de acceso---, se pueden separar y decodificar
cada una por separado, permitiendo que todos los dispositivos avancen en el proceso de
conexión sin usar recursos adicionales~\cite{ref4, ref5}. La base técnica de este enfoque es
el acceso múltiple no ortogonal (NOMA, \textit{Non-Orthogonal Multiple Access}), que al
aplicarse al contexto del acceso aleatorio da origen a NORA.

NOMA se basa en dos mecanismos que trabajan juntos: la codificación por superposición en el
transmisor y la cancelación sucesiva de interferencia (SIC,
\textit{Successive Interference Cancellation}) en el receptor~\cite{ref5, ref8}. En el
enlace ascendente, varios dispositivos transmiten al mismo tiempo sobre el mismo recurso
tiempo-frecuencia pero con diferentes niveles de potencia. El eNB recibe la suma
de estas señales y aplica SIC para separarlas: primero decodifica la señal más fuerte
tratando a las demás como ruido, después la resta de la señal total y decodifica la
siguiente, y así sucesivamente hasta recuperar cada mensaje. Para que SIC funcione bien, es
necesario que las señales tengan diferencias de potencia claras, lo que en el acceso aleatorio
se consigue gracias a la variación natural del canal entre dispositivos que están a diferentes
distancias del eNB~\cite{ref5}.

En NORA, la separación de señales se hace desde el primer mensaje del procedimiento de
acceso. Cuando dos o más UEs transmiten el mismo preámbulo y sus señales llegan con tiempos
de arribo distintos ---porque están a diferentes distancias del eNB---, el
receptor puede estimar el ToA de cada dispositivo y asignarle un nivel de potencia
diferente~\cite{ref4}. Con esta información, el eNB genera varias respuestas RAR
para el mismo preámbulo colisionado, una por cada dispositivo identificado, y le pide a cada
uno que ajuste su potencia de transmisión antes de enviar MSG3. Cuando recibe los MSG3, el
eNB aplica SIC para separar las transmisiones superpuestas y decodificar cada
mensaje de forma independiente~\cite{ref4, ref5}.

El impacto de esta modificación en el desempeño es importante. Liang
\textit{et al.}~\cite{ref4} demuestran que NORA supera al esquema ortogonal convencional en
probabilidad de éxito del acceso, probabilidad de colisión y capacidad de la red, con una
mejora de más del 30\,\% en el número de dispositivos que se pueden atender al mismo tiempo.
Por su parte, Seo \textit{et al.}~\cite{ref5} muestran que el \textit{throughput} máximo con
NORA puede superar 0{,}7 paquetes por ranura, frente al límite teórico de 0{,}368 del
esquema ALOHA ranurado\footnote{ALOHA es un protocolo de acceso múltiple por contención en
el que los dispositivos transmiten sin coordinación previa; en su versión ranurada
(\textit{Slotted ALOHA}) las transmisiones se alinean a ranuras de tiempo, lo que reduce las
colisiones y eleva el \textit{throughput} máximo respecto al ALOHA puro.} convencional. Wang \textit{et al.}~\cite{ref8} amplían este análisis
a redes MTC masivas, donde la formación de grupos NOMA y la optimización del factor de
retroceso de potencia permiten superar al NORA convencional en \textit{throughput} cuando
la carga es alta.

A partir del esquema básico se han propuesto varias variantes para escenarios específicos.
El acceso aleatorio no ortogonal sin concesión (GF-NORA,
\textit{Grant-Free Non-Orthogonal Random Access}) elimina la fase de solicitud de recursos y
permite que los dispositivos transmitan directamente sus datos sin esperar una asignación
previa, lo que reduce la latencia pero requiere mayor esfuerzo para detectar la actividad
de los usuarios~\cite{ref9}. El NORA triangular (T-NORA, \textit{Triangular NORA}) adapta el
procedimiento a paquetes de distinto tamaño en sistemas de antenas masivas~\cite{ref11},
mientras que el NORA con paginación grupal (NORA-GP, \textit{NORA-based Group Paging})
controla la sobrecarga en redes 5G con tráfico MTC intenso~\cite{ref12}. En el caso
particular de NB-IoT, Wang \textit{et al.}~\cite{ref7} proponen el NORA para NB-IoT
(NB-NORA), que combina la identificación por ToA con control de potencia adaptativo y logra un
\textit{throughput} hasta un 75\,\% mayor que el del esquema ortogonal de referencia.

Sin embargo, todas estas variantes comparten una misma limitación: separar señales por
potencia requiere que los niveles recibidos sean suficientemente diferentes y que el número
de dispositivos colisionados no sea demasiado grande para que SIC funcione bien. Cuando se
quiere ir más allá de lo que permite la separación por potencia, es necesario agregar otra
forma de distinguir a los dispositivos. Esto lleva al uso de técnicas basadas en códigos
escasos, como SCMA, que permiten diferenciar dispositivos no solo por su potencia sino
también por el patrón de recursos de frecuencia que usan, como se describe en la siguiente
sección.


% ------------------------------------------------------------
\section{Multiplexación por código escaso y separación multiusuario}
% ------------------------------------------------------------

La cancelación sucesiva de interferencia descrita en la sección anterior funciona mejor
cuando las señales de los dispositivos colisionados tienen diferencias claras de potencia en
el receptor. Sin embargo, esto no siempre ocurre: dispositivos que están a distancias
parecidas del eNB generan niveles de potencia similares, lo que dificulta que
el receptor SIC pueda separarlos correctamente~\cite{ref10}. Para superar esta limitación,
el acceso múltiple por código escaso (SCMA, \textit{Sparse Code Multiple Access}) agrega una
forma adicional de separación basada en códigos, que funciona sin depender de la potencia
recibida y es especialmente útil en escenarios de acceso masivo con transmisión no coordinada.

SCMA pertenece a la familia de técnicas de acceso múltiple no ortogonal en el dominio del
código de baja densidad (LDCD-NOMA,
\textit{Low-Density Code-Domain NOMA})~\cite{ref10, ref17}. Su idea principal es asignar a
cada dispositivo un \textit{codebook}, que corresponde a una matriz compleja dispersa. Ese
\textit{codebook} define a qué subportadoras del enlace ascendente se asignan los símbolos de
datos del usuario. Lo clave de esa asignación es su escasez (\textit{sparsity}), ya que cada
\textit{codebook} usa solo unas pocas de las subportadoras disponibles y deja las demás en
cero. Gracias a esto, varios dispositivos con \textit{codebooks} diferentes pueden transmitir
al mismo tiempo sobre las mismas subportadoras sin que sus señales sean iguales, porque cada
uno usa un patrón de subportadoras distinto~\cite{ref10, ref17}. Por ejemplo, si hay cuatro
subportadoras y cada \textit{codebook} usa solo dos, se pueden superponer hasta seis usuarios
sobre los mismos recursos de frecuencia con patrones distintos, lo que aumenta de forma
notable la capacidad del sistema comparado con el acceso ortogonal.

Para separar las señales superpuestas en el receptor se usa el algoritmo de paso de
mensajes (MPA, \textit{Message Passing Algorithm}), que trabaja sobre una representación
del sistema como un grafo bipartito\footnote{Un grafo bipartito es aquel cuyos nodos se
dividen en dos conjuntos disjuntos, de modo que toda arista conecta un nodo de un conjunto
con uno del otro y nunca dos nodos del mismo conjunto.}~\cite{ref10, ref17}. En ese grafo
hay dos tipos de nodos: los nodos variables, que representan los dispositivos activos, y los
nodos función, que representan las subportadoras compartidas. Una arista conecta un
dispositivo con una subportadora cuando su \textit{codebook} la usa. Sobre esa estructura,
el MPA intercambia mensajes de forma iterativa entre ambos conjuntos de nodos y actualiza
las estimaciones de probabilidad sobre los símbolos transmitidos hasta llegar a una decisión.
Lo importante es que, como el grafo es disperso ---cada dispositivo está conectado a pocas
subportadoras y cada subportadora es compartida por pocos dispositivos---, la complejidad
computacional se reduce significativamente en comparación con la de un receptor óptimo de
máxima verosimilitud~\cite{ref17}. A diferencia del SIC, el MPA no necesita que las señales tengan
niveles de potencia diferentes, lo que lo hace más robusto cuando los dispositivos están a
distancias parecidas del eNB~\cite{ref10}.

En este trabajo, SCMA se aplica al procedimiento de acceso aleatorio de NB-IoT para ampliar
la capacidad de separación de dispositivos colisionados más allá de lo que permite la
diferencia de ToA o de potencia por sí sola. En el esquema CRACS-R
(\cref{sec:cracs_r}), el eNB responde a cada preámbulo colisionado
con dos RAR, cada uno con seis \textit{codebooks} disponibles, de manera que los dispositivos
colisionados pueden elegir combinaciones distintas de RAR y \textit{codebook} para transmitir
su MSG3~\cite{ref10}. El detector MPA en recepción procesa los MSG3 recibidos sobre los
recursos compartidos e intenta recuperar cada mensaje por separado, avanzando en el proceso
de acceso solo a los dispositivos cuya separación sea exitosa.

La combinación de \textit{codebooks} escasos y detección MPA hace que SCMA sea una herramienta muy
útil para aumentar la capacidad del canal de acceso aleatorio: permite que más dispositivos
compartan los mismos recursos físicos sin necesitar más ancho de banda ni diferencias de
potencia entre ellos. Sin embargo, el desempeño de este esquema también depende de las
condiciones del canal. Cuando la transmisión se hace a través de un enlace satelital de
órbita baja, el canal introduce efectos adicionales ---retardo variable, desplazamiento
Doppler y desvanecimiento dinámico--- que no existen en redes terrestres y que afectan
directamente la detección MPA y el procedimiento de acceso en general. Antes de abordar el
entorno satelital, es necesario describir cómo el eNB detecta la ocurrencia de una
colisión, ya que esto determina directamente qué tan bien funcionan los esquemas NORA y SCMA.


% ------------------------------------------------------------
\section{Detección de colisiones y métricas de clasificación}
\label{sec:deteccion_colisiones}
% ------------------------------------------------------------

Un aspecto clave en los esquemas NORA y SCMA es que el eNB pueda
detectar la ocurrencia de una colisión en MSG1, antes de enviar la respuesta RAR. La
efectividad de esta detección afecta directamente el desempeño del sistema, ya que una
colisión que no se detecta produce un solo RAR para varios dispositivos y, como
consecuencia, una falla en MSG3 que obliga a repetir todo el procedimiento.

La detección de colisiones se basa en analizar el perfil de retardo de potencia (PDP,
\textit{Power Delay Profile}), que se obtiene al correlacionar la señal recibida en el PRACH
con la secuencia Zadoff-Chu generada en el eNB~\cite{ref20}. Cuando un solo
dispositivo transmite un preámbulo, el PDP muestra un solo pico de correlación. Cuando varios
dispositivos transmiten el mismo preámbulo desde diferentes distancias, aparecen varios picos
dentro de la misma zona de correlación cero (ZCZ, \textit{Zero Correlation Zone}), lo que
indica que hubo una colisión~\cite{ref20}.

Estudios recientes han mostrado que las técnicas de aprendizaje automático pueden mejorar
mucho la precisión de esta detección. Maldonado Cardenas \textit{et al.}~\cite{ref20}
proponen un método basado en redes neuronales que, entrenado con los 24 valores de potencia
del PDP por compartimiento, alcanza una precisión balanceada (\textit{Balanced Accuracy})
mayor al 98\,\% en escenarios con condiciones de canal homogéneas y del 95\,\% en
escenarios con condiciones de canal cruzadas (es decir, cuando se entrena y se prueba con
modelos de canal diferentes).

Para medir qué tan bien un sistema detecta colisiones, se usa la matriz de confusión, una
tabla asociada a las predicciones del detector que las clasifica en cuatro categorías, como
se muestra en la Tabla~\ref{tab:confusion_matrix}.

\begin{table}[h!]
\centering
\caption{Estructura de la matriz de confusión para detección de colisiones}
\label{tab:confusion_matrix}
\begin{tabular}{|l|c|c|}
\hline
 & \textbf{Colisión real} & \textbf{No colisión real} \\
\hline
\textbf{Detectada como colisión} & Verdadero Positivo (VP) & Falso Positivo (FP) \\
\textbf{Detectada como no colisión} & Falso Negativo (FN) & Verdadero Negativo (VN) \\
\hline
\end{tabular}
\end{table}

A partir de esta matriz se calculan las métricas que describen el desempeño del detector:

\begin{itemize}
    \item \textbf{Precisión} (\textit{Precision}): fracción de detecciones de colisión que
          corresponden efectivamente a colisiones reales, que se estima de acuerdo con la
          Ecuación~\ref{eq:precision}.
          \begin{equation}
          \text{Precisión} = \frac{VP}{VP + FP}
          \label{eq:precision}
          \end{equation}

    \item \textbf{Exhaustividad} (\textit{Recall}): fracción de colisiones reales que el
          detector identifica correctamente, según la Ecuación~\ref{eq:recall}.
          \begin{equation}
          \text{Recall} = \frac{VP}{VP + FN}
          \label{eq:recall}
          \end{equation}

    \item \textbf{Especificidad} (\textit{Specificity}): fracción de no-colisiones que el
          detector clasifica bien, de acuerdo con la Ecuación~\ref{eq:especificidad}. Esta
          métrica es muy importante porque una especificidad baja significa que se están
          clasificando como colisiones preámbulos que en realidad fueron exitosos, lo que
          desperdicia recursos que se habrían usado correctamente~\cite{ref20}.
          \begin{equation}
          \text{Especificidad} = \frac{VN}{VN + FP}
          \label{eq:especificidad}
          \end{equation}

    \item \textbf{Precisión balanceada} (\textit{Balanced Accuracy}): promedio entre
          \textit{Recall} y Especificidad, calculado con la Ecuación~\ref{eq:bacc}, que da una
          evaluación equilibrada del desempeño en ambas clases. Es especialmente útil cuando
          hay un desbalance entre la cantidad de colisiones y no-colisiones.
          \begin{equation}
          \text{Balanced Accuracy} = \frac{\text{Recall} + \text{Especificidad}}{2}
          \label{eq:bacc}
          \end{equation}
\end{itemize}

En el esquema CRACS-R, el eNB usa una probabilidad de detección de colisión para
decidir si envía dos RAR (cuando detecta colisión) o uno solo (cuando no la detecta). Esta
probabilidad suele ser alta, cercana a 0{,}975. Conceptualmente, la matriz de
confusión admite dos tipos de error:

\begin{itemize}
    \item Un \textbf{falso negativo} (colisión no detectada) implica que los dispositivos
          colisionados no reciben la respuesta diferenciada: el eNB emite un único RAR,
          todos los UEs colisionados transmiten MSG3 sobre los mismos recursos sin
          separación por \textit{codebook}, ninguno se puede decodificar y los UEs deben
          reintentar al expirar el \textit{timer} de contención de MSG3. Esto incrementa
          el retardo de acceso y reduce la capacidad efectiva.

    \item Un \textbf{falso positivo} (no-colisión clasificada como colisión) lleva al
          eNB a generar dos RAR innecesarios, desperdiciando recursos de
          señalización y potencialmente confundiendo al UE que transmitió sin colisión.
\end{itemize}

En este trabajo el modelo de simulación considera únicamente la probabilidad
de detección de colisión (sensibilidad) y, con ella, la tasa de falsos negativos. Los
falsos positivos se dejan fuera del modelo: representan un preámbulo sin colisión
tratado como colisionado y su inclusión implicaría simular señalización adicional sin
usuarios reales detrás, añadiendo complejidad sin contribuir a la pregunta de
investigación sobre cuánto baja la capacidad del canal al degradarse la detección.


% ------------------------------------------------------------
\section{Redes satelitales de órbita baja y entornos no terrestres (NTN)}
% ------------------------------------------------------------

En las secciones anteriores se describieron las tecnologías y técnicas de acceso consideradas
en este trabajo asumiendo que el canal entre el dispositivo y el eNB es
relativamente estable y que los retardos de propagación son muy pequeños, como pasa en las
redes celulares terrestres. Sin embargo, cuando la infraestructura de acceso se mueve al
espacio, estas condiciones cambian por completo. Las redes no terrestres (NTN,
\textit{Non-Terrestrial Networks}) ---que incluyen satélites de diferentes órbitas y
plataformas de gran altitud--- se han vuelto un complemento necesario de las redes celulares
para llevar conectividad a regiones donde la infraestructura terrestre no alcanza o no
existe~\cite{ref3, ref13}. Entre las distintas órbitas posibles, los satélites de órbita
baja (LEO, \textit{Low Earth Orbit}) son la opción más prometedora para aplicaciones IoT de
baja latencia, gracias a que están relativamente cerca de la superficie
terrestre~\cite{ref3, ref6}.

Los satélites LEO orbitan entre los 500 y 2000 km de altitud, lo que los diferencia de los
satélites de órbita media (MEO, \textit{Medium Earth Orbit}) y de los geoestacionarios
(GEO, \textit{Geostationary Earth Orbit}), que están a unos
35\,786 km~\cite{ref6, ref18}. Esta menor altitud ofrece dos ventajas principales: una
pérdida de trayectoria (\textit{path loss}) mucho menor que en GEO y un retardo de ida y
vuelta (RTT, \textit{Round-Trip Time}) de entre 2 y 20 ms según la elevación, comparado con
los aproximadamente 500 ms de la órbita geoestacionaria~\cite{ref6}. Sin embargo, la baja
altitud también implica una velocidad orbital alta ---alrededor de 7{,}5 km/s---, lo que
significa que cada satélite es visible desde un punto fijo en tierra durante solo cinco a
quince minutos, dependiendo de la altitud y la elevación mínima~\cite{ref18}. Esto tiene
una consecuencia directa sobre el acceso: todos los dispositivos dentro del haz del satélite
que necesiten conectarse deben hacerlo en esa breve ventana de visibilidad, lo que genera
ráfagas de intentos de acceso simultáneo que agravan el problema de colisiones comparado
con las redes terrestres~\cite{ref13}.

Además, el movimiento del satélite introduce efectos en el canal que no existen en los
sistemas terrestres fijos. El más importante es el desplazamiento Doppler: el movimiento
relativo entre el satélite y los dispositivos en tierra produce un cambio en la frecuencia
portadora recibida que, en la banda Ku, puede superar los 200 kHz en el peor
caso~\cite{ref6}. Este desplazamiento afecta la ortogonalidad entre subportadoras en sistemas
OFDM y dificulta la detección de preámbulos en NPRACH si no se compensa adecuadamente.
Aunque se pueden usar técnicas basadas en sistemas globales de navegación por satélite (GNSS,
\textit{Global Navigation Satellite System}) para estimar y precompensar parcialmente el
Doppler, los desfases que quedan ---causados por multitrayecto y errores en la estimación---
pueden llegar a varios kilohertz y deben manejarse en el receptor~\cite{ref6}. También se debe
considerar que el retardo de propagación cambia continuamente a medida que el satélite
se mueve, lo cual afecta directamente el cálculo del avance de temporización en NB-IoT, un
parámetro crítico para sincronizar el enlace ascendente. El canal satelital presenta además
un componente dominante de línea de vista (LOS, \textit{Line of Sight}) en la mayoría de los
casos, que se modela con distribuciones de Rician, aunque en zonas suburbanas o con
obstáculos puede haber condiciones de no visibilidad directa (NLOS, \textit{Non-Line of
Sight})~\cite{ref13}.

Ante estos desafíos, 3GPP ha incorporado soporte para NTN en sus especificaciones desde la
Release 16, con una integración más completa en la Release 17, donde se formaliza la
arquitectura de acceso radio no terrestre tanto para NR (\textit{New Radio}) como para
NB-IoT~\cite{ref13}. La Release 18 identifica la congestión en el acceso como un problema
abierto en escenarios NTN-IoT y plantea soluciones para versiones futuras~\cite{ref13}. Esta
evolución en las normas muestra que simplemente extender los procedimientos terrestres al
entorno satelital no es suficiente: las condiciones cambiantes del canal LEO, la variación
del retardo y las ráfagas de acceso hacen necesario adaptar los mecanismos de acceso
aleatorio a este entorno.

Es en este contexto donde se sitúa el presente trabajo. La combinación de NB-IoT con enlaces
satelitales LEO es un escenario relevante pero todavía poco estudiado desde el punto de vista
de la capacidad del canal de acceso no ortogonal~\cite{ref4, ref6, ref13}. El comportamiento
de esquemas como NORA y SCMA bajo las condiciones de un enlace LEO ---retardo variable,
Doppler residual, canal Rician--- no se puede deducir directamente de los resultados
obtenidos en redes terrestres, lo que justifica el estudio por simulación que se plantea en
este trabajo. La siguiente sección del marco teórico presenta las métricas que
permiten medir la capacidad del canal de acceso en este escenario.


% ------------------------------------------------------------
\section{Capacidad del canal de acceso aleatorio}
% ------------------------------------------------------------

El concepto de capacidad en el acceso aleatorio es diferente de la definición clásica de
Shannon, que establece la tasa máxima a la que se puede transmitir información con una
probabilidad de error muy pequeña sobre un canal dado~\cite{ref16}. Mientras que la capacidad
de Shannon se refiere al límite teórico del enlace físico entre un transmisor y un receptor,
la capacidad del canal de acceso aleatorio es la medida que mejor se ajusta a este tipo de
escenarios: indica la cantidad máxima de dispositivos que pueden completar el procedimiento
de acceso de forma simultánea, con los recursos disponibles y bajo las condiciones de
operación del sistema~\cite{ref4, ref5, ref10}. Es la métrica adecuada porque en escenarios
MTC el cuello de botella generalmente no es la tasa del enlace físico sino la eficiencia del
mecanismo de acceso cuando muchos dispositivos compiten por los mismos recursos.

La capacidad del canal de acceso se evalúa usando varias métricas que miden diferentes
aspectos del desempeño del sistema. La probabilidad de éxito de acceso ($P_s$)
indica la fracción de dispositivos que logran completar la conexión en un intento dado, sin
necesitar retransmisión; su complemento, la probabilidad de fallo, depende de las colisiones
de preámbulo y las fallas al decodificar MSG3~\cite{ref4, ref10}. El
\textit{throughput} de acceso aleatorio ($T$) es el número de intentos de acceso
exitosos por unidad de tiempo o por ranura de acceso, y es la métrica que mejor refleja qué
tan bien se aprovecha el canal; en el acceso ortogonal clásico tipo ALOHA ranurado (S-ALOHA,
\textit{Slotted ALOHA}), el \textit{throughput} máximo teórico es $1/e \approx 0{,}368$
paquetes por ranura, un límite que los esquemas NORA y SCMA buscan superar~\cite{ref5}. El
retardo de acceso ($D$) mide el tiempo que pasa desde que un dispositivo inicia el
procedimiento de acceso hasta que lo completa, incluyendo todos los reintentos; en NB-IoT,
este retardo puede ser considerable incluso sin colisiones debido a las periodicidades del
NPRACH, y se vuelve especialmente crítico en los niveles CE1 y CE2~\cite{ref7}. Por último, la
eficiencia de acceso relaciona el número de dispositivos aceptados con el total de
recursos consumidos por el procedimiento, tomando en cuenta tanto el canal de preámbulo como
el canal de datos ascendente.

Para este trabajo, la métrica principal es la cantidad de usuarios aceptados de forma
simultánea, es decir, el número máximo de dispositivos que pueden completar el procedimiento de acceso
en una ventana de tiempo dada sin que el sistema se sature~\cite{ref4}. Esta cantidad depende
de tres factores relacionados entre sí: la densidad de dispositivos activos, qué tan bien el
esquema de acceso resuelve las colisiones, y las condiciones del canal. En el entorno LEO,
el tercer factor cobra mayor importancia: el retardo variable, el Doppler residual y las
ráfagas de acceso que se producen durante la ventana de visibilidad del satélite cambian la
distribución de los intentos de acceso en el tiempo y afectan la probabilidad de decodificar
MSG3 correctamente, lo que reduce la capacidad efectiva respecto a lo que predicen los
modelos terrestres~\cite{ref6, ref13}.

La relación entre la densidad de dispositivos y la capacidad del canal se puede ver en la
curva de \textit{throughput} del sistema. A medida que el número de dispositivos activos $M$
crece, el \textit{throughput} aumenta al principio porque hay más transmisiones útiles, pero
a partir de cierto punto las colisiones empiezan a dominar y el \textit{throughput} comienza
a bajar, hasta que el sistema colapsa por congestión~\cite{ref5}. El punto de máximo
\textit{throughput} define la carga óptima de operación y, junto con la probabilidad de éxito
de acceso bajo esa carga, marca la capacidad operativa del canal. Un esquema como NORA, que
recupera dispositivos colisionados en vez de descartarlos, mueve este punto de máximo hacia
cargas mayores, aumentando la capacidad del canal sin necesitar recursos
adicionales~\cite{ref4, ref5}. Estas métricas son las que se usan para comparar los esquemas
Legacy, CRACS-R y CRACS-T en la simulación de este trabajo.
